% Teil II: der durchgerechnete Weg.

% Gedankenschritt: Quantor umdrehen + Prädikat verneinen; De Morgan für
% Verknüpfungen; Wahrheitswert spielt keine Rolle.

\subsection*{Gegeben und gesucht}
Zu vier Aussagen das Gegenteil (die Negation) angeben, ohne zu bewerten,
ob die Aussage wahr ist.

\subsection*{Ansatz}
\textbf{Regel: Quantoren umdrehen, Prädikat verneinen, ``und``/``oder``
vertauschen.} Dabei wird der Wahrheitsgehalt nicht beurteilt, sondern nur
die logische Struktur ``gespiegelt``.

\subsection*{a) ``Die Sonne geht jeden Morgen auf.''}
\textbf{Struktur:} $\forall m \in \text{Morgen}: \text{Sonne geht auf}$.
Negation: $\exists m: \neg(\text{geht auf})$.
Gegenteil: \emph{Es gibt einen Morgen, an dem die Sonne nicht aufgeht.}

\subsection*{b) ``Jedes Tier ist grau.''}
$\forall$ Tier $t$: $t$ ist grau. Negation: $\exists t$: $t$ ist nicht
grau.
Gegenteil: \emph{Es gibt ein Tier, das nicht grau ist.}

\subsection*{c) ``Es gibt eine Pflanze, die weder grün noch gelb ist.''}
\textbf{Struktur:} $\exists p$: $\neg(\text{grün } \text{oder}\ \text{gelb})$.
Negation (Existenzquantor zu Allquantor, Prädikat verneint; De Morgan):
$\forall p$: $\neg\neg(\text{grün oder gelb})$
$=\ \forall p$: grün $=$ oder gelb.
Gegenteil: \emph{Jede Pflanze ist grün oder gelb.}

\subsection*{d) ``Jede Bitfolge besteht aus Nullen und Einsen.''}
$\forall$ Bitfolge $b$: $b$ besteht aus Nullen und Einsen. Negation:
$\exists b$: $b$ besteht nicht nur aus Nullen und Einsen.
Gegenteil: \emph{Es gibt eine Bitfolge, die nicht nur aus Nullen und
Einsen besteht.}

\begin{ergebnis}
\begin{itemize}
  \item a) Es gibt einen Morgen, an dem die Sonne nicht aufgeht.
  \item b) Es gibt ein Tier, das nicht grau ist.
  \item c) Jede Pflanze ist grün oder gelb.
  \item d) Es gibt eine Bitfolge, die nicht nur aus Nullen und Einsen besteht.
\end{itemize}
\end{ergebnis}

\begin{kontrolle}
\textbf{Logische Probe:} Aussage und Negation müssen sich gegenseitig
ausschliessen und zusammen alle Fälle überdecken. Für a) stehen
``jeden Morgen geht die Sonne auf`` und ``mindestens ein Morgen ohne
Sonnenaufgang`` als Gegensätze da; beide zugleich wahr wären ein
Widerspruch, also ist die Negation richtig.

\textbf{Typischer Fehler (ausgeschlossen):} ``Kein Tier ist grau`` wäre
viel zu stark. Gäbe es einige graue und einige nicht graue Tiere, wäre
weder ``jedes Tier ist grau`` noch ``kein Tier ist grau`` wahr, und die
beiden würden sich nicht korrekt ergänzen --- die Negation
``mindestens eines ist nicht grau`` ist die korrekte.
\end{kontrolle}

\begin{merksatz}
Negation ``dreht`` Quantoren: ``alle`` $\leftrightarrow$ ``mindestens
eine``, ``es gibt`` $\leftrightarrow$ ``keine``; De Morgan für
``und``/``oder``; Wahrheitswert wird nicht bestimmt.
\end{merksatz}
